\section{Introduction}
\label{sec:intro}

Two paradigms now dominate robot learning. \emph{Vision-language-action} (VLA)
models predict actions from frozen weights; \emph{code-as-policy} agents instead
write executable programs and improve them from experience~\citep{codeaspolicies,aspire}.
This survey is organized around that fault line.

\paragraph{Why now: the skill stack.}
The app store for robot skills has arrived---Unitree UniStore ships one-tap,
cross-model motion downloads~\citep{unistore2026}. But every skill in it is
\emph{canned playback}: the least-capable point on our map. The skill stack has
three layers---\textbf{(1) Build} (how skills are created, \S\ref{sec:taxonomy}--\S\ref{sec:techniques}),
\textbf{(2) Distribute} (marketplaces and hubs), and \textbf{(3) Adapt} (on-device
self-improvement, \S\ref{sec:code}). UniStore is layer~2 over static layer-1 skills;
the promised value---``adjust grip per fruit so nothing bruises''---lives in layer~3,
which no shipped skill does today.

\paragraph{Contributions.}
(i) a taxonomy of robot-learning techniques organized by \emph{weights vs.\ skills}
(\S\ref{sec:taxonomy}); (ii) a deep-dive arranging code-as-policy by \emph{degree of
self-improvement} and isolating the un-surveyed \S3.1.5 cell (\S\ref{sec:code});
(iii) a companion map of the ``skills'' pole (\S\ref{sec:skill}); and (iv) an
open-problems agenda for the emerging skill economy (\S\ref{sec:open}).

\paragraph{Related surveys.}
Table~\ref{tab:survey_compare} contrasts the topical coverage of the closest surveys with
ours, and Table~\ref{tab:surveys} places them on a timeline. Within \emph{ACM Computing
Surveys} the nearest are on world models, embodied intelligence, and language-model
navigation, and beyond it a cluster covers vision-language-action models and foundation models
for manipulation---but none organises code-as-policy by degree of self-improvement or names
the \S3.1.5 cell.
% Table — coverage of related surveys vs. ours.
\begin{table}[t]
\centering
\footnotesize
\renewcommand{\arraystretch}{1.25}
\begin{tabular}{@{}p{0.24\textwidth}c c *{6}{c}@{}}
\toprule
\textbf{Survey} & \textbf{Venue} & \textbf{Yr}
  & \tweights{} & \tcode{} & \tskill{} & \treward{}
  & \textcolor{secE!55!black}{\textbf{Self-impr.}} & \textcolor{cat-bench!45!black}{\textbf{Economy}}\\
\midrule
World Models~\cite{survey_worldmodels}          & CSUR  & 2025 & \pp & \nn & \nn & \nn & \nn & \nn\\
Embodied Intelligence~\cite{survey_embodiedintel}& CSUR  & 2025 & \pp & \nn & \pp & \nn & \nn & \nn\\
Navigation via FLM~\cite{survey_navigation}      & CSUR  & 2026 & \pp & \pp & \nn & \nn & \nn & \nn\\
VLA for Embodied AI~\cite{survey_vla}            & arXiv & 2024 & \yy & \nn & \nn & \nn & \nn & \nn\\
Large-VLM VLA~\cite{survey_vlmvla}               & arXiv & 2025 & \yy & \nn & \nn & \pp & \nn & \nn\\
Embodied Learning~\cite{survey_embodiedlearn}    & MIR   & 2025 & \pp & \nn & \pp & \pp & \nn & \nn\\
FM for Manipulation~\cite{survey_fmrobot}        & arXiv & 2024 & \yy & \pp & \pp & \pp & \nn & \nn\\
\midrule
\textbf{Ours (\emph{Weights or Skills?})}        & ---   & 2026 & \yy & \yy & \yy & \yy & \yy & \yy\\
\bottomrule
\end{tabular}
\caption{Topical coverage of the closest related surveys versus ours.
\yy\ covered, \pp\ partial/mentioned, \nn\ not covered. Columns: \tweights{} (VLA/weights),
\tcode{} (code-as-policy), \tskill{} (skill libraries), \treward{} (reward synthesis),
\textbf{Self-impr.}\ (organised by \emph{degree of self-improvement}, incl.\ the \S3.1.5 cell),
\textbf{Economy} (the robot skill marketplace). Only this survey covers the last two axes.
CSUR = ACM Computing Surveys; MIR = Machine Intelligence Research.}
\label{tab:survey_compare}
\end{table}

% Table 1 — related surveys as a colored TikZ timeline (year axis, one marker each).
\begin{table}[t]
\centering
\resizebox{\columnwidth}{!}{%
\begin{tikzpicture}[font=\small]
  % year gridlines + labels
  \foreach \x/\yr in {0/2023, 2.4/2024, 4.8/2025, 7.2/2026}{
    \draw[gray!22] (\x,-0.35) -- (\x,5.5);
    \node[gray, font=\footnotesize, below] at (\x,-0.4) {\yr};
  }
  % rows: dotted leader + name (left) + colored marker at its year
  \draw[gray!35, densely dotted] (-0.25,5) -- (0,5);
  \node[anchor=east] at (-0.3,5) {\textbf{Foundation Models in Robotics}};
  \filldraw[fill=secA, draw=hidden-draw] (0,5) circle (5.5pt);

  \draw[gray!35, densely dotted] (-0.25,4) -- (2.4,4);
  \node[anchor=east] at (-0.3,4) {\textbf{VLA Models survey}};
  \filldraw[fill=secB, draw=hidden-draw] (2.4,4) circle (5.5pt);

  \draw[gray!35, densely dotted] (-0.25,3) -- (2.4,3);
  \node[anchor=east] at (-0.3,3) {\textbf{LLMs for Robotics}};
  \filldraw[fill=secC, draw=hidden-draw] (2.4,3) circle (5.5pt);

  \draw[gray!35, densely dotted] (-0.25,2) -- (2.4,2);
  \node[anchor=east] at (-0.3,2) {\textbf{Robot Skill Learning}};
  \filldraw[fill=secF, draw=hidden-draw] (2.4,2) circle (5.5pt);

  \draw[gray!35, densely dotted] (-0.25,1) -- (4.8,1);
  \node[anchor=east] at (-0.3,1) {\textbf{Code Generation for Robotics}};
  \filldraw[fill=secD, draw=hidden-draw] (4.8,1) circle (5.5pt);

  \draw[gray!35, densely dotted] (-0.25,0) -- (7.2,0);
  \node[anchor=east] at (-0.3,0) {\textbf{Ours --- \emph{Weights or Skills?}}};
  \node[star, star points=5, star point ratio=2.3, fill=secE, draw=hidden-draw,
        minimum size=18pt, inner sep=0] at (7.2,0) {};
\end{tikzpicture}%
}
\caption{Related surveys on a timeline. Prior surveys (2023--2025) each cover part of the
landscape---foundation models, VLA policies, LLM planning, skill acquisition, code
generation---but none organizes code-as-policy by \emph{degree of self-improvement}.
\textbf{Ours} ($\star$, 2026) isolates the feedback\,+\,memory\,+\,search cell (\S3.1.5)
that ASPIRE occupies.}
\label{tab:surveys}
\end{table}

\FloatBarrier
